%% LyX 2.4.4 created this file.  For more info, see https://www.lyx.org/.
%% Do not edit unless you really know what you are doing.
\documentclass[english]{article}
\usepackage[T1]{fontenc}
\usepackage[latin9]{inputenc}
\usepackage{enumitem}
\usepackage{amsmath}
\usepackage{amsthm}
\usepackage{amssymb}
\usepackage{geometry}
\geometry{verbose,tmargin=1in,bmargin=1in,lmargin=1in,rmargin=1in}

\makeatletter
%%%%%%%%%%%%%%%%%%%%%%%%%%%%%% Textclass specific LaTeX commands.
\numberwithin{equation}{section}
\numberwithin{figure}{section}
\newlength{\lyxlabelwidth}      % auxiliary length 

%%%%%%%%%%%%%%%%%%%%%%%%%%%%%% User specified LaTeX commands.
\usepackage{fancyhdr}
\usepackage{lastpage}
\pagestyle{fancy}
\usepackage{enumitem}
\usepackage{ifthen}
\everymath=\expandafter{\the\everymath\displaystyle}
%\DeclareMathSizes{10}{10}{10}{10}
\usepackage{multicol}

\usepackage{tikz}
\usetikzlibrary{patterns}
\usetikzlibrary{plotmarks}
\usepackage{pgfplots}
\pgfplotsset{compat=newest}

\usepackage{calc}
\usepackage[nomessages]{fp}

\usepackage{datetime}


%%%%%commands
\newcommand{\mb}[1]{\mathbb{#1}}
\newcommand{\funct}[3]{$#1\colon#2\rightarrow#3$}
% Added by lyx2lyx
\setlength{\parskip}{\smallskipamount}
\setlength{\parindent}{0pt}

\makeatother

\usepackage{babel}
\begin{document}
\renewcommand{\author}{Dr.\,James Ashe}

\newcommand{\classNum}{439}

\newcommand{\classSec}{A}

\newcommand{\nameType}{Name: \hspace{-4cm}}

\newcommand{\examType}{Final}

\newcommand{\Date}{\formatdate{1}{5}{2014}} %%\AdvanceDate[12] \today

%\newdate{my_date}{\day}{\month}{\year}%   
%\AdvanceDate[-5] 
%\displaydate{my_date}

%%First page's header%%%%%%%%%%%%%%%%%%%%%%
\fancypagestyle{firststyle} %%header settings for first pagr
{
	\fancyhead[L]{MTH \classNum{}(\classSec{}) \author{}\\
	\textbf{\examType{}} \Date{}}
	\fancyhead[C]{\textbf{\nameType}}
	\setlength{\headsep}{14pt}
}
%%%%%%%%%%%%%%%%%%%%%%%%%%%%%%%%%%%%%%%%%%

%%Next page's header%%%%%%%%%%%%%%%%%%%%%%%
%\setlist{nolistsep} %eliminates space between list items
\setlist[enumerate]{topsep=-2pt,itemsep=6pt}
\lhead{MTH \classNum{}(\classSec{})} 
\chead{\textbf{\nameType}} 
\cfoot{\thepage{}~of~\pageref{LastPage}} 
\setlength{\headsep}{5pt} 
%%%%%%%%%%%%%%%%%%%%%%%%%%%%%%%%%%%%%%%%%%%

%%Various settings; see the preamble for other settings%%%%%
%%\setenumerate[0]{leftmargin=12pt,itemindent=0pt,labelwidth=0pt}
\renewcommand{\labelenumi}{\textbf{\arabic{enumi}.}}
\renewcommand{\labelenumii}{\textbf{\alph{enumii}.}}
\thispagestyle{firststyle}
%%%%%%%%%%%%%%%%%%%%%%%%%%%%%%%%%%%%%%%%%%
\newcommand{\xspace}{\vspace{.85cm}}

\subsubsection*{Rings and Subrings. \textmd{Choose any 3.}}
\begin{enumerate}[resume]
\item Rings: \emph{Give an example for each of the following.}

\begin{enumerate}
\item A ring with identity
\item A ring without identity
\item A non-commutative ring
\item A subring of $\mb{Z}_{6}$, $\mb{Q}$, and $\mb{C}$.
\end{enumerate}
\item Prove that the set of nonnegative integers is not a subring of $\mathbb{Z}$.
\item Prove that $\mathbb{Z}[i]=\left\{ a+bi\colon a,b\in\mathbb{Z}\right\} $
is a subring of $\mathbb{C}$.
\item Let $R$ be a ring with subrings $A$ and $B$. Prove that $A\cap B$
is also a subring. \emph{\vfill}
\end{enumerate}

\subsubsection*{Integral Domains and Fields. \textmd{Choose any 3.}}
\begin{enumerate}[resume]
\item Integral domains and Fields: \emph{Give an example for each of the
following.}

\begin{enumerate}
\item An integral domain.
\item A ring which is not an integral domain.
\item A field.
\item A ring which is not a field. 
\end{enumerate}
\item Let $R$ be an integral domain and suppose that $ax-bx=0$. Prove
that if $a\neq b$ then $x=0$.
\item Let $R=\mathbb{Z}_{10}$ 

\begin{enumerate}
\item Prove that $R$ is not an integral domain. 
\item Prove that $R$ is not a field.
\end{enumerate}
\item Let $R=\mathbb{Z}_{3}$

\begin{enumerate}
\item Prove that $R$ is a field.
\item Prove that $R$ is an integral domain.\emph{\vfill}
\end{enumerate}
\end{enumerate}

\subsubsection*{Ideals. \textmd{Choose any 3.}}
\begin{enumerate}[resume]
\item Ideals: \emph{Give an example for each of the following.}

\begin{enumerate}
\item A ring and an ideal in that ring.
\item A ring and a subring which is not an ideal.
\item A ring other than $\mathbb{Z}$, and a principal ideal in that ring.
\item A ring other than $\mathbb{Z}_{12}$, and a prime ideal in that ring.
\end{enumerate}
\item Prove that for any ring $R$, $0_{R}$ is an ideal of $R$.
\item Prove that $E=\{x\colon x=2n\mbox{ for }n\in\mathbb{Z}\}$ is an ideal
in $\mathbb{Z}$.
\item Find all the ideals and maximal ideals of $\mathbb{Z}_{12}$ (use
a lattice). \emph{\vfill}
\end{enumerate}

\subsubsection*{Homomorphisms and Isomorphisms. \textmd{Choose any 3.}}
\begin{enumerate}[resume]
\item Homomorphism and Isomorphisms: \emph{Give an example for each of the
following.}

\begin{enumerate}
\item A function between rings which is not 1-1.
\item An onto function between rings.
\item An isomorphism between rings (other than the identity function).
\item A homomorphism between rings which is not an isomorphism.
\end{enumerate}
\item Prove that the rings $\mb{Q}$ and $\mb{Z}$ are not isomorphic.
\item Prove that $\mb{Z}$ is not isomorphic to $\mb{Z}_{6}$.
\item Let $f\colon R\rightarrow S$ be a ring homomorphism.

\begin{enumerate}
\item What is $f(0_{R})$?
\item Suppose that $f(a)f(b)=1_{R}$ find $f(ab)$.
\item Suppose that $\ker f=\left\{ 0,r\right\} $ for some nonzero $r\in R$.
Is $R$ isomorphic to $S$?\vfill
\end{enumerate}
\end{enumerate}

\subsubsection*{Factor Rings. \textmd{Choose any 3.}}
\begin{enumerate}[resume]
\item Ideals: \emph{Give an example for each of the following}

\begin{enumerate}
\item A factor ring.
\item Two elements in the factor ring above.
\item An abelian factor ring. 
\item A factor ring that is a field. 
\end{enumerate}
\item Let $F=\left\{ x\colon x=4n+1\mbox{ for }n\in\mathbb{Z}\right\} $,
the set of odd integers. 

\begin{enumerate}
\item Prove that $F$ is not an ideal in $\mathbb{Z}$.
\item Prove that $\mathbb{Z}/F$ is not a ring. 
\end{enumerate}
\item Consider $\mathbb{Z}_{10}/\left\langle 2\right\rangle $. 

\begin{enumerate}
\item Give two examples of distinct elements in $\mathbb{Z}_{10}/\left\langle 2\right\rangle $.
\item Prove that $\mathbb{Z}_{10}/\left\langle 2\right\rangle $ is an integral
domain
\end{enumerate}
\item Define $f:\mathbb{Z}\rightarrow\mathbb{Z}_{6}$ as $f(x)=x\bmod6$.

\begin{enumerate}
\item Let $f:\mathbb{Z}\rightarrow\mathbb{Z}_{6}$. Find the kernel of $f$.
\item Show that $\mathbb{Z}/6\mathbb{Z}$ is isomorphic to $\mathbb{Z}_{6}$.
\end{enumerate}
\item Recall that $\left\langle i\right\rangle $ is a principal ideal in
$\mathbb{Z}[i]$. Suppose that the function $f\colon\mathbb{Z}[i]\rightarrow\mathbb{Z}$
is an onto homomorphism with kernel, $\ker f=\left\langle i\right\rangle $
. Prove that $\mathbb{Z}[i]/\left\langle i\right\rangle $ is isomorphic
to $\mathbb{Z}$.\emph{\vfill}
\end{enumerate}

\end{document}
